\documentclass[12pt]{article}
\usepackage{notes-common}
\usetikzlibrary{calc}
\geometry{margin=0.35in}
\AtBeginDocument{\fontsize{14}{17}\selectfont}

% Mon Oct 5 -- Ohm's law + Kirchhoff's laws (Knight 27.5, 28.1-28.2). Same
% lesson for both tracks (no calculus needed), so no PHYS 230/235 split --
% matches electric-potential-notes.tex's layout conventions (small margins,
% bigger base font). See circuits-intro-notes-key.tex for the filled-in
% version -- same structure, wording and diagrams, every blank carries the
% answer. Built from Seth's rough notes (uploaded).
\newcommand{\eqline}[2][-0.5mm]{%
  \par\vspace{1mm}\noindent #2\par\vspace{#1}%
  \noindent\textcolor{rulehue}{\rule{\linewidth}{0.5pt}}\par\vspace{4mm}%
}

% ---- simple battery/resistor schematic pieces, shared by all three circuit
% diagrams below. A loop is drawn corner-to-corner with \draw, then a
% battery or resistor symbol is dropped on top of the segment it sits on.
\newcommand{\battv}[1]{% vertical battery symbol centered at (#1)
  % explicit leads, guaranteeing the plates connect to the rest of the
  % loop even if the background wire doesn't render cleanly behind them
  \draw[line width=0.7pt] ($(#1)+(0,0.55)$) -- ($(#1)+(0,0.12)$);
  \draw[line width=0.7pt] ($(#1)+(0,-0.12)$) -- ($(#1)+(0,-0.55)$);
  \draw[line width=1.1pt] ($(#1)+(-0.35,0.12)$) -- ($(#1)+(0.35,0.12)$);
  \draw[line width=2.4pt] ($(#1)+(-0.2,-0.12)$) -- ($(#1)+(0.2,-0.12)$);
  \node[scale=0.75] at ($(#1)+(0.6,0.12)$) {$+$};
  \node[scale=0.75] at ($(#1)+(0.6,-0.12)$) {$-$};
}
\newcommand{\reszig}[1]{% horizontal resistor zigzag centered at (#1)
  \fill[white] ($(#1)+(-0.65,-0.32)$) rectangle ($(#1)+(0.65,0.32)$);
  \draw ($(#1)+(-0.6,0)$)
    -- ($(#1)+(-0.42,0.22)$) -- ($(#1)+(-0.24,-0.22)$)
    -- ($(#1)+(-0.06,0.22)$) -- ($(#1)+(0.12,-0.22)$)
    -- ($(#1)+(0.3,0.22)$) -- ($(#1)+(0.48,-0.22)$)
    -- ($(#1)+(0.6,0)$);
}

\begin{document}

\noteshead{Battery \& Resistor Circuits}{}

\goal{
\begin{itemize}[itemsep=5pt,parsep=2pt,topsep=2pt]
\item Pin down what current and $\Delta V$ actually mean, and the four
rules that follow just from charge and energy conservation.
\item Turn Ohm's law into a tool: given a resistor and a $\Delta V$ across
it, get $I$.
\item Work a single resistor both with an ideal battery and with a real
one (internal resistance $r$), to see what $r$ costs you.
\item Derive $R_S=R_1+R_2$ for series and $R_P=(1/R_1+1/R_2)^{-1}$ for
parallel from the rules in \S1 --- not as formulas to memorize, but as
consequences of ``same current'' and ``same $\Delta V$.''
\end{itemize}}

%============================================================
\nsec{Definitions}

\textbf{Current}: \bl[8cm]

\hspace{1cm} units: \bl[2.5cm] $= $ \bl[2.5cm]

\textbf{Change in voltage, $\Delta V$}: \bl[9cm]

\vspace{2mm}
\textbf{Series}: elements connected \bl[5cm]

\textbf{Parallel}: elements connected \bl[5cm]

%============================================================
\nsec{Four rules}

\begin{itemize}[itemsep=7pt,parsep=3pt,topsep=3pt]
\item Elements in series have \bl[6cm]
\item Elements in parallel have \bl[6cm]
\item Current is \bl[4cm] at a junction.
\item Around any closed loop, the $\Delta V$'s \bl[5cm]
\end{itemize}

\checkpoint{The last two rules are each a conservation law in disguise. Which
one is charge conservation, and which one is energy conservation?}

%============================================================
\nsec{Ohm's law defines a resistor}

\[
\Delta V = \
\]
\vspace{0.3cm}

Given a resistor of resistance $R$ and a voltage difference $\Delta V$
across it, the current through it is

\eqline{$I = $}

%============================================================
\nsec{One resistor and a battery}

\begin{center}
\begin{tikzpicture}[x=1cm,y=1cm]
  \draw (0,0) -- (0,2) -- (3,2) -- (3,0) -- cycle;
  \battv{0,1}
  \node[scale=0.8, anchor=east] at (-0.55,1) {$\varepsilon,\,r$};
  \reszig{1.5,2}
  \node[scale=0.8, anchor=south] at (1.5,2.45) {$R$};
\end{tikzpicture}
\end{center}

\textbf{(a) Ideal battery}: $\varepsilon=12$ V, $r=0$, $R=4\ \Omega$.

\eqline[0.1cm]{$I = $}
\eqline{$V_B = $}

\textbf{(b) Real battery}: same $\varepsilon$ and $R$, but now $r=1\ \Omega$.

\eqline[0.1cm]{$I = $}
\eqline{$V_B = $}

\checkpoint{$V_B$ dropped in part (b). Where did that energy per charge go?}

%============================================================
\nsec{Two resistors in series}

\begin{center}
\begin{tikzpicture}[x=1cm,y=1cm]
  \draw (0,0) -- (0,2) -- (6,2) -- (6,0) -- cycle;
  \battv{0,1}
  \node[scale=0.8, anchor=east] at (-0.55,1) {$\varepsilon$};
  \reszig{2,2}
  \node[scale=0.8, anchor=south] at (2,2.45) {$R_1$};
  \reszig{4,2}
  \node[scale=0.8, anchor=south] at (4,2.45) {$R_2$};
\end{tikzpicture}
\end{center}

\textbf{Derive $R_S$} (symbols only --- use the series current rule and Ohm's
law on each resistor, then the loop rule):

\derivbox{4.3cm}{start from $\varepsilon = \Delta V_1 + \Delta V_2$}

\result[1.1cm]{$R_S = $}

\textbf{Now with numbers}: $\varepsilon=12$ V, $R_1=4\ \Omega$, $R_2=8\ \Omega$
(ideal battery).

\eqline[0.1cm]{$R_S = $}
\eqline[0.1cm]{$I = $}
\eqline[0.1cm]{$\Delta V_1 = $ \hspace{3.2cm} $\Delta V_2 = $}

\checkpoint{Check the loop rule with your two $\Delta V$'s above: do they
add up to $\varepsilon$?}

%============================================================
\nsec{Two resistors in parallel}

\begin{center}
\begin{tikzpicture}[x=1cm,y=1cm]
  \draw (0,0) -- (0,2) -- (1.3,2) -- (1.3,2.9) -- (4.7,2.9) -- (4.7,2) -- (6,2) -- (6,0) -- cycle;
  \draw (1.3,2) -- (1.3,1.1) -- (4.7,1.1) -- (4.7,2);
  \battv{0,1}
  \node[scale=0.8, anchor=east] at (-0.55,1) {$\varepsilon$};
  \reszig{3,2.9}
  \node[scale=0.8, anchor=south] at (3,3.35) {$R_1$};
  \reszig{3,1.1}
  \node[scale=0.8, anchor=north] at (3,0.65) {$R_2$};
\end{tikzpicture}
\end{center}

\textbf{Derive $R_P$} (symbols only --- use the parallel voltage rule and
Ohm's law on each resistor, then the junction rule at the top node):

\derivbox{4.3cm}{start from $I = I_1 + I_2$}

\result[1.1cm]{$R_P = $}

\textbf{Now with numbers}: $\varepsilon=12$ V, $R_1=4\ \Omega$,
$R_2=12\ \Omega$ (ideal battery).

\eqline[0.1cm]{$R_P = $}
\eqline[0.1cm]{$I_1 = $ \hspace{3.2cm} $I_2 = $}
\eqline[0.1cm]{$I_{\text{total}} = $}

\checkpoint{Which resistor carries more current, and why does that make
sense given which one is smaller?}

\end{document}
